% !TeX root = ../main.tex
% =================================================================================
\chapter{Introduction}\label{chap:introduction}
% =================================================================================
% The introduction summarizes your entire thesis. A reader must understand the context, the
% problem, your goal and research questions, your approach, your results, and your contributions
% from the introduction alone. Write it last, but collect notes for it from the beginning.
% The introduction usually has no sections, except for the unnumbered "Thesis Structure" at the end.

% --- Context -----------------------------------------------------------------------
Start with the context of your thesis.
Introduce the general area your work belongs to and explain why it is relevant, ideally with concrete examples from practice or research.
Introduce every term with a specific meaning in precise language that a fellow computer science student understands.
Support your statements with references, e.g., to foundational work~\cite{Lamport1994}.

% --- Problem -----------------------------------------------------------------------
Next, describe the problem you address.
Explain the limitations of existing approaches or tools and why they are not sufficient.
Be concrete: Who is affected by the problem, and what are the consequences?
At the end of these paragraphs, the reader should understand the gap between the current state and a desired state.

% --- Goal and research questions ---------------------------------------------------
Then, state the goal of your thesis and how it addresses this gap.
Derive your research questions from the problem and state each of them where it naturally arises in your argument.
For example: Our goal is to enable programs to cooperate safely when they modify the same files.
To this end, a program first has to notice changes made by others, which leads to our first research question:
%
\defRQ{How can a program detect that another program modified a file it saves?}{detect}
%
Once a modification is detected, the user has to decide which changes to keep.
We investigate how this decision can be supported, leading to our second research question:
%
\defRQ{How can users resolve conflicting modifications of a file?}{resolve}
%
Refer to research questions with \verb|\refRQ{label}|, e.g., \refRQ{detect}.

% --- Approach and results ----------------------------------------------------------
Briefly describe how you approach the research questions, e.g., which methodology you follow (\cref{chap:methodology}), how you elicit requirements, what you design and implement, and how you evaluate it.
Then summarize your main results with concrete numbers, e.g., \enquote{In a user study with 12~participants, our approach achieves a System Usability Scale score of 81.2, compared to 31.2 for the baseline.}
If you publish a replication package, mention it here.

% --- Benefits and limitations ------------------------------------------------------
State who benefits from your results (e.g., developers, end users, researchers) and how.
Also name the most important limitations that remain for future work.

% --- Contributions -----------------------------------------------------------------
In summary, our core contributions are as follows:
\begin{contributions}
  \item A concept for \ldots\label{c:concept}
  \item An open-source implementation of \ldots, available at \url{https://github.com/...}.\label{c:implementation}
  \item An evaluation of \ldots through a user study with \ldots participants.\label{c:evaluation}
\end{contributions}
Refer to contributions with \verb|\refC{label}|, e.g., \refC{concept}, just like research questions.

\section*{Thesis Structure}

The remainder of this thesis is structured as follows:
\begin{description}
  \item[\Cref{chap:foundations} -- \nameref{chap:foundations}:] We introduce the concepts and technologies required for understanding this thesis.
  \item[\Cref{chap:methodology} -- \nameref{chap:methodology}:] We describe the research methodology used in this thesis.
  \item[\Cref{chap:related-work} -- \nameref{chap:related-work}:] We give an overview of related work and delimit our work from it.
  \item[\Cref{chap:requirements} -- \nameref{chap:requirements}:] We describe the requirements for our approach and how we elicited them.
  \item[\Cref{chap:concept} -- \nameref{chap:concept}:] We present our concept for \ldots
  \item[\Cref{chap:implementation} -- \nameref{chap:implementation}:] We describe the architecture and selected implementation details.
  \item[\Cref{chap:evaluation} -- \nameref{chap:evaluation}:] We evaluate our approach and answer the research questions.
  \item[\Cref{chap:conclusion} -- \nameref{chap:conclusion}:] We conclude the thesis and outline future work.
\end{description}
